% --- CORRECCIÓN IMPORTANTE EN LA PRIMERA LÍNEA ---
% Agregamos "xcolor=table" para que funcione \rowcolor sin errores
\documentclass[aspectratio=169, 10pt, xcolor=table]{beamer}

% --- Configuración del Tema (estilo profesional como el ejemplo) ---
\usetheme{Madrid}
\usecolortheme{dolphin}

% --- Paquetes necesarios ---
\usepackage[utf8]{inputenc}
\usepackage[spanish]{babel}
\usepackage{amsmath, amssymb}
\usepackage{booktabs}
\usepackage{graphicx}
\usepackage{listings}
\usepackage{tikz}
\usetikzlibrary{arrows.meta, positioning, shapes.geometric}

% --- Ruta de Imágenes (crea la carpeta "imagenes" y coloca ahí tus archivos) ---
\graphicspath{{./imagenes/}}

% --- Estilo para código Python ---
\definecolor{codegreen}{rgb}{0,0.6,0}
\definecolor{codegray}{rgb}{0.5,0.5,0.5}
\definecolor{codepurple}{rgb}{0.58,0,0.82}
\definecolor{backcolour}{rgb}{0.95,0.95,0.92}
\lstset{
    language=Python,
    backgroundcolor=\color{backcolour},
    commentstyle=\color{codegreen},
    keywordstyle=\color{blue},
    numberstyle=\tiny\color{codegray},
    stringstyle=\color{codepurple},
    basicstyle=\ttfamily\scriptsize,
    breaklines=true,
    frame=single,
    showstringspaces=false
}

% --- Pie de página personalizado ---
\makeatletter

\setbeamertemplate{footline}{
  \leavevmode%
  \hbox{%
  \begin{beamercolorbox}[wd=.333333\paperwidth,ht=2.25ex,dp=1ex,center]{author in head/foot}%
    \usebeamerfont{author in head/foot}\insertshortauthor
  \end{beamercolorbox}%
  \begin{beamercolorbox}[wd=.333333\paperwidth,ht=2.25ex,dp=1ex,center]{title in head/foot}%
    \usebeamerfont{title in head/foot}Sesión 7
  \end{beamercolorbox}%
  \begin{beamercolorbox}[wd=.333333\paperwidth,ht=2.25ex,dp=1ex,right]{date in head/foot}%
    \usebeamerfont{date in head/foot}NLP \hfill \insertframenumber/\inserttotalframenumber \hspace*{0.3cm}
  \end{beamercolorbox}}%
  \vskip0pt%
}

\makeatother

% --- Datos de la portada ---
\title[SESIÓN 7 NLP]{\textbf{Sesión 7}}
\subtitle{Redes Recurrentes (RNN, LSTM, GRU) para NLP}
\author[Prof. C. Sánchez]{Carlos César Sánchez Coronel}
\institute{\textbf{NLP}}
\date{}

% Logo opcional (descomentar si tienes la imagen)
% \titlegraphic{\includegraphics[height=1.5cm]{logo.png}}

% --- Eliminar la barra de navegación (opcional) ---
\setbeamertemplate{navigation symbols}{}

% --- Índice al inicio de cada sección ---
\AtBeginSection[]{
  \begin{frame}
    \frametitle{Contenido}
    \tableofcontents[currentsection]
  \end{frame}
}

% ============================================================================
% INICIO DEL DOCUMENTO
% ============================================================================
\begin{document}

% --- Portada ---
\begin{frame}
    \titlepage
\end{frame}

% --- Índice general ---
\begin{frame}{Contenido}
    \tableofcontents
\end{frame}

% ============================================================================
% SECCIÓN 1: INTRODUCCIÓN
% ============================================================================
\section{Introducción}

\begin{frame}{La necesidad de modelar secuencias}
    \begin{itemize}
        \item Las redes feed-forward (con average pooling) pierden el orden de las palabras.
        \item No distinguen: ``El producto es bueno, pero la atención es pésima'' vs ``La atención es pésima, pero el producto es bueno''.
        \item El lenguaje es inherentemente secuencial: el significado depende del contexto previo.
    \end{itemize}
    % Basado en Pregunta 1
\end{frame}

% ============================================================================
% SECCIÓN 2: RNN VANILLA
% ============================================================================
\section{RNN Vanilla}

\begin{frame}{Arquitectura de una RNN}
    \begin{block}{Actualización del estado oculto}
        \[
        h_t = \tanh(W_{hh} h_{t-1} + W_{xh} x_t + b_h)
        \]
        \[
        y_t = \text{softmax}(W_{hy} h_t + b_y)
        \]
        \begin{itemize}
            \item $x_t$: entrada en el paso $t$ (embedding de palabra).
            \item $h_t$: estado oculto actual (memoria de la secuencia).
        \end{itemize}
    \end{block}
    % IMAGEN: Diagrama de RNN desenrollada
\end{frame}

\begin{frame}{Backpropagation Through Time (BPTT)}
    \begin{itemize}
        \item Desenrollar la RNN para cada paso de tiempo.
        \item La pérdida total es la suma de pérdidas en cada paso.
        \item El gradiente se propaga hacia atrás a través de todos los pasos.
        \item Complejidad: $O(T)$ en tiempo, pero secuencial.
    \end{itemize}
    % Basado en Pregunta 5
\end{frame}

\begin{frame}{Problema del gradiente desvaneciente/explosivo}
    \begin{itemize}
        \item Durante BPTT, el gradiente se multiplica por $W_{hh}$ repetidamente.
        \item Si $\|W_{hh}\| < 1$, el gradiente se desvanece (tiende a cero).
        \item Si $\|W_{hh}\| > 1$, el gradiente explota.
        \item \textbf{Consecuencia}: la RNN no puede aprender dependencias de largo alcance (ej. una palabra al inicio de un párrafo que afecta una palabra al final).
    \end{itemize}
    % Basado en Pregunta 2
\end{frame}

% ============================================================================
% SECCIÓN 3: LSTM (LONG SHORT-TERM MEMORY)
% ============================================================================
\section{LSTM (Long Short-Term Memory)}

\begin{frame}{Arquitectura LSTM}
    \begin{block}{Componentes}
        \begin{itemize}
            \item \textbf{Puerta de olvido} ($f_t$): qué información del estado de celda anterior olvidar.
            \item \textbf{Puerta de entrada} ($i_t$): qué nueva información almacenar.
            \item \textbf{Candidato a memoria} ($\tilde{c}_t$): nueva información propuesta.
            \item \textbf{Puerta de salida} ($o_t$): qué parte del estado de celda emitir.
        \end{itemize}
    \end{block}
    % Basado en Pregunta 3
\end{frame}

\begin{frame}{Ecuaciones LSTM}
    \begin{align*}
        f_t &= \sigma(W_f[h_{t-1}, x_t] + b_f) \\
        i_t &= \sigma(W_i[h_{t-1}, x_t] + b_i) \\
        \tilde{c}_t &= \tanh(W_c[h_{t-1}, x_t] + b_c) \\
        c_t &= f_t \odot c_{t-1} + i_t \odot \tilde{c}_t \\
        o_t &= \sigma(W_o[h_{t-1}, x_t] + b_o) \\
        h_t &= o_t \odot \tanh(c_t)
    \end{align*}
    % Basado en Pregunta 3
\end{frame}

\begin{frame}{Interpretación de las puertas}
    \begin{itemize}
        \item \textbf{Puerta de olvido} ($f_t$): valores cercanos a 0 = olvidar, 1 = recordar. Permite descartar información irrelevante.
        \item \textbf{Puerta de entrada} ($i_t$): controla cuánto del nuevo candidato se añade.
        \item \textbf{Estado de celda} ($c_t$): autopista lineal para el gradiente (sumas y productos), evitando el desvanecimiento.
        \item \textbf{Puerta de salida} ($o_t$): decide qué exponer al resto de la red.
    \end{itemize}
    % Basado en Pregunta 12
\end{frame}

\begin{frame}{Ventajas de LSTM}
    \begin{itemize}
        \item Maneja dependencias a largo plazo (cientos de pasos).
        \item El gradiente fluye fácilmente a través de $c_t$ debido a operaciones lineales.
        \item Fue el estado del arte en muchas tareas secuenciales antes de los Transformers.
    \end{itemize}
\end{frame}

% ============================================================================
% SECCIÓN 4: GRU (GATED RECURRENT UNIT)
% ============================================================================
\section{GRU (Gated Recurrent Unit)}

\begin{frame}{Arquitectura GRU}
    \begin{block}{Componentes}
        \begin{itemize}
            \item \textbf{Puerta de reinicio} ($r_t$): cuánto del estado anterior olvidar.
            \item \textbf{Puerta de actualización} ($z_t$): combina olvido y entrada de LSTM.
        \end{itemize}
    \end{block}
    % Basado en Pregunta 15
\end{frame}

\begin{frame}{Ecuaciones GRU}
    \begin{align*}
        r_t &= \sigma(W_r[h_{t-1}, x_t] + b_r) \\
        z_t &= \sigma(W_z[h_{t-1}, x_t] + b_z) \\
        \tilde{h}_t &= \tanh(W_h[r_t \odot h_{t-1}, x_t] + b_h) \\
        h_t &= (1 - z_t) \odot h_{t-1} + z_t \odot \tilde{h}_t
    \end{align*}
    % Basado en Pregunta 15
\end{frame}

\begin{frame}{Comparación LSTM vs GRU}
    \begin{table}
        \centering
        \begin{tabular}{l|c|c}
            \toprule
            \textbf{Característica} & \textbf{LSTM} & \textbf{GRU} \\
            \midrule
            Número de puertas & 3 (forget, input, output) & 2 (update, reset) \\
            Estado de celda separado & Sí & No (solo estado oculto) \\
            Parámetros & Más & Menos \\
            Velocidad de entrenamiento & Más lenta & Más rápida \\
            Riesgo de sobreajuste & Mayor (con pocos datos) & Menor \\
            \bottomrule
        \end{tabular}
    \end{table}
    % Basado en Pregunta 4
\end{frame}

% ============================================================================
% SECCIÓN 5: RNN BIDIRECCIONALES
% ============================================================================
\section{RNN Bidireccionales}

\begin{frame}{RNN Bidireccional}
    \begin{itemize}
        \item Dos RNN: una procesa de izquierda a derecha (forward), otra de derecha a izquierda (backward).
        \item Estados concatenados: $h_t = [\overrightarrow{h_t}; \overleftarrow{h_t}]$.
        \item Ideal para tareas donde se tiene acceso a toda la secuencia (NER, clasificación de oraciones).
    \end{itemize}
    % Basado en Pregunta 6
\end{frame}

% ============================================================================
% SECCIÓN 6: APLICACIONES EN NLP
% ============================================================================
\section{Aplicaciones en NLP}

\begin{frame}{Clasificación de sentimiento}
    \begin{itemize}
        \item Usar el último estado oculto o promedio de todos los estados.
        \item LSTM bidireccional + capa fully-connected + softmax.
    \end{itemize}
    % Basado en Pregunta 10
\end{frame}

\begin{frame}{Generación de texto}
    \begin{itemize}
        \item Modelo de lenguaje: predecir la siguiente palabra.
        \item Entrenamiento con cross-entropy.
        \item Generación: greedy search o sampling con temperatura.
        \item Evaluación: perplejidad (perplexity).
    \end{itemize}
    % Basado en Pregunta 8
\end{frame}

\begin{frame}{Traducción automática (seq2seq)}
    \begin{itemize}
        \item Encoder RNN: codifica frase origen en un vector de contexto.
        \item Decoder RNN: genera frase destino a partir del contexto.
        \item Limitación: vector de contexto fijo para frases largas.
        \item Mejora: mecanismo de atención (próxima sesión).
    \end{itemize}
    % Basado en Pregunta 16
\end{frame}

\begin{frame}{Reconocimiento de Entidades Nombradas (NER)}
    \begin{itemize}
        \item Etiquetado por token (B-PER, I-PER, O, etc.).
        \item LSTM bidireccional + capa lineal + softmax en cada paso.
        \item Pérdida: cross-entropy por token.
    \end{itemize}
    % Basado en Pregunta 13
\end{frame}

% ============================================================================
% SECCIÓN 7: EJEMPLO EN PYTORCH (CLASIFICACIÓN DE SENTIMIENTO)
% ============================================================================
\section{Ejemplo en PyTorch: Clasificación de Sentimiento con LSTM}

\begin{frame}[fragile]{Datos y vocabulario}
    \begin{lstlisting}
import torch
import torch.nn as nn
import torch.optim as optim
from collections import Counter

texts = [
    "I loved this movie, it was absolutely fantastic",
    "Terrible film, I hated it, worst movie ever",
    "Amazing acting and great story",
    "Boring and slow, waste of time"
]
labels = [1, 0, 1, 0]

def build_vocab(texts, max_size=10000):
    word_counts = Counter()
    for text in texts:
        word_counts.update(text.lower().split())
    vocab = {word: idx+2 for idx, (word, _) in enumerate(word_counts.most_common(max_size))}
    vocab['<PAD>'] = 0
    vocab['<UNK>'] = 1
    return vocab

vocab = build_vocab(texts)
vocab_size = len(vocab)
    \end{lstlisting}
\end{frame}

\begin{frame}[fragile]{Codificar y definir modelo}
    \begin{lstlisting}
def encode(text, vocab, max_len=20):
    tokens = text.lower().split()[:max_len]
    indices = [vocab.get(token, vocab['<UNK>']) for token in tokens]
    indices += [vocab['<PAD>']] * (max_len - len(indices))
    return indices

X = torch.tensor([encode(t, vocab) for t in texts])
y = torch.tensor(labels)

class LSTMSentiment(nn.Module):
    def __init__(self, vocab_size, embed_dim=50, hidden_dim=64, num_classes=2):
        super().__init__()
        self.embedding = nn.Embedding(vocab_size, embed_dim, padding_idx=0)
        self.lstm = nn.LSTM(embed_dim, hidden_dim, batch_first=True)
        self.fc = nn.Linear(hidden_dim, num_classes)
        self.dropout = nn.Dropout(0.5)

    def forward(self, x):
        embedded = self.embedding(x)
        out, (hidden, cell) = self.lstm(embedded)
        last_hidden = hidden[-1]
        out = self.fc(self.dropout(last_hidden))
        return out
    \end{lstlisting}
\end{frame}

\begin{frame}[fragile]{Entrenamiento y evaluación}
    \begin{lstlisting}
model = LSTMSentiment(vocab_size)
criterion = nn.CrossEntropyLoss()
optimizer = optim.Adam(model.parameters(), lr=0.001, weight_decay=1e-5)

for epoch in range(200):
    optimizer.zero_grad()
    logits = model(X)
    loss = criterion(logits, y)
    loss.backward()
    optimizer.step()
    if (epoch+1) % 50 == 0:
        print(f'Epoch {epoch+1}, Loss: {loss.item():.4f}')

with torch.no_grad():
    logits = model(X)
    preds = torch.argmax(logits, dim=1)
    acc = (preds == y).float().mean()
    print(f'Accuracy: {acc:.2%}')
    \end{lstlisting}
\end{frame}

% ============================================================================
% SECCIÓN 8: MÉTRICAS
% ============================================================================
\section{Métricas en NLP con RNN}

\begin{frame}{Perplejidad (Perplexity)}
    \[
    PPL = \exp\left(-\frac{1}{N}\sum_{i=1}^{N} \log P(w_i | \text{contexto})\right)
    \]
    \begin{itemize}
        \item Mide qué tan ``sorprendido'' está el modelo por los datos.
        \item Menor perplejidad = mejor modelo.
        \item Límite inferior: 1 (modelo perfecto); superior: tamaño del vocabulario (uniforme).
    \end{itemize}
    % Basado en Pregunta 9
\end{frame}

\begin{frame}{Otras métricas}
    \begin{itemize}
        \item \textbf{Accuracy}: para clasificación de secuencias completas.
        \item \textbf{Precisión, recall, F1-score}: para tareas de etiquetado (NER, POS).
        \item \textbf{Exactitud de secuencia}: coincidencia exacta de toda la secuencia generada.
    \end{itemize}
\end{frame}

% ============================================================================
% SECCIÓN 9: LIMITACIONES DE RNN/LSTM
% ============================================================================
\section{Limitaciones de RNN/LSTM}

\begin{frame}{Limitaciones}
    \begin{itemize}
        \item \textbf{Procesamiento secuencial}: no paralelizable a lo largo del tiempo, lento en GPU.
        \item \textbf{Dependencias muy largas}: aunque LSTM mejora, aún puede tener dificultades con secuencias de miles de pasos.
        \item \textbf{Memoria limitada}: el estado oculto tiene capacidad finita.
    \end{itemize}
    % Basado en Pregunta 11
\end{frame}

% ============================================================================
% SECCIÓN 10: PREPARACIÓN PARA ENTREVISTAS
% ============================================================================
\section{Preparación para Entrevistas}

\begin{frame}{Preguntas típicas}
    \begin{itemize}
        \item ¿Por qué LSTM funciona mejor que RNN vanilla para secuencias largas?
        \item Diferencia entre LSTM y GRU.
        \item ¿Qué problema resuelve la puerta de olvido en LSTM?
        \item ¿Cómo se entrena una RNN? (BPTT)
        \item ¿Cuándo usarías una RNN en lugar de una feed-forward?
    \end{itemize}
\end{frame}

\begin{frame}{Preguntas avanzadas}
    \begin{itemize}
        \item Explica el truncamiento en BPTT.
        \item ¿Cómo manejarías secuencias de longitud variable en un lote? (padding, pack\_padded\_sequence)
        \item ¿Qué son las RNN bidireccionales?
        \item ¿Cómo inicializarías los pesos de las puertas en LSTM?
        \item ¿Qué es el dropout variacional en RNN?
    \end{itemize}
    % Basado en preguntas 5, 7, 18, 19
\end{frame}

% ============================================================================
% SECCIÓN 11: CONEXIÓN CON EL RESTO DEL CURSO
% ============================================================================
\section{Conexión con el resto del curso}

\begin{frame}{Próximos pasos}
    \begin{itemize}
        \item \textbf{Sesión 4 (Word Embeddings)}: base para capas de embedding.
        \item \textbf{Sesión 5 (Modelos Clásicos)}: superados por RNN en tareas secuenciales.
        \item \textbf{Sesión 8 (Atención y Transformers)}: resuelven las limitaciones de RNN (paralelización, dependencias largas).
    \end{itemize}
\end{frame}

% ============================================================================
% SECCIÓN 12: RESUMEN
% ============================================================================
\section{Resumen}

\begin{frame}{Lo que hemos aprendido}
    \begin{exampleblock}{Conceptos clave}
        \begin{itemize}
            \item \textbf{RNN vanilla}: procesan secuencias pero sufren gradiente desvaneciente.
            \item \textbf{LSTM}: puertas de olvido, entrada, salida; estado de celda preserva gradientes.
            \item \textbf{GRU}: simplificación de LSTM, menos parámetros, comparable rendimiento.
            \item \textbf{BPTT}: algoritmo de entrenamiento para RNN.
            \item \textbf{RNN bidireccionales}: contexto completo, útiles para NER.
            \item \textbf{Aplicaciones}: clasificación, generación, traducción, NER.
            \item \textbf{Métricas}: perplejidad, F1, accuracy.
            \item \textbf{Limitaciones}: secuencialidad, memoria finita.
        \end{itemize}
    \end{exampleblock}
\end{frame}

% --- Diapositiva final ---
\begin{frame}
    \centering
    \Huge \textbf{¡Gracias!}

\end{frame}

\end{document}